\documentclass[12pt]{article}

\usepackage{graphicx}
\usepackage{hyperref}
\usepackage{float}

\title{PPC Práctica 2 - UDP, JSON y XML}
\date{14 de Noviembre, 2026}
\author{Ángel Ruiz Fernandez}

\begin{document}

    \maketitle
    \pagebreak

    \tableofcontents
    \pagebreak

    \section{Funcionalidades}

    Se han implementado todos los requisitos obligatorios:

    \begin{itemize}
    \item a. mensajes de distribución con cabecera de \texttt{Content-Type} generados
    \item b. mensajes de control
    \begin{itemize}
        \item cambiar codificación
        \item cambiar frencuencia
        \item cambiar unidades
        \item activar y desactivar envío
    \end{itemize}
    \item formatos XML y JSON:
    \begin{itemize}
        \item validación XML Schema Definition 1.0 para XML utilizando \texttt{javax.xml} con archivos \texttt{.xsd} incluidos
        \item parser DOM de XML
        \item JSON utilizando paquete GSON
        \item volcado a un archivo \texttt{server.log} en el servidor
    \end{itemize}
    \item requisitos opcionales
    \begin{itemize}
        \item distintas tramas de control
        \item distintos clientes en misma maquina (gracias a \texttt{setsockopt(SO\_REUSEADDR)})
    \end{itemize}
    \end{itemize}

    Mensajes descritos en \texttt{MENSAJES.md}

    Ya que el servidor ha de enviar y recibir datos simultaneamente, es natural implementarlo
    mediante dos hilos.

    La validación mediante XML Schema Definition es la más extendida y con mayor funcionalidad, sencilla de usar con \texttt{javax.xml}.

    El grupo multicast utilizado es configurable en tiempo de compilación modificando la constante \texttt{group}, que por defecto es \texttt{ff12::cafe}.
    Asimismo, el puerto configirable de destino de datagramas es el \texttt{6767}, y el de comandos de control el \texttt{6768}.

    \section{Implementación}
    
    \subsection{Diseño}

    Utilizando sockets multicast IPv6 con MulticastSocket bindeados a una interfaz especifica link-local.

    Cliente y servidor usan E/S síncrona, ya que al tener un thread dedicados a escuchar o enviar,
    pueden esperar síncronamente si es hora de enviar o si se han recibido datos nuevos.

    La serialización y deserialización es ad-hoc, ya que es más sencilla y compacta.
    La serialización/deserialización mediante Java Reflection requiere una clase Bean en GSON
    y en XML es inexistente porque XML no corresponde 1:1 a objetos, mientras que JSON si.

    \subsubsection{Formato de datagramas}

    \begin{verbatim}
    Content-Type: <tipo>

    [datos en xml o json]
    \end{verbatim}

    Tipo puede ser según el formato de la trama

    \begin{itemize}
    \item \texttt{application/xml}
    \item \texttt{application/json}
    \end{itemize}

    Datagrama XML:
    \begin{verbatim}
<?xml version="1.0"?>
<xs:schema xmlns:xs="http://www.w3.org/2001/XMLSchema">

  <xs:element name="data">
    <xs:complexType>
      <xs:sequence>
        <xs:element name="server" type="xs:string"/>

        <xs:choice minOccurs="0" maxOccurs="unbounded">

          <xs:element name="measurement">
            <xs:complexType>
              <xs:simpleContent>
                <xs:extension base="xs:decimal">
                  <xs:attribute name="magnitude" use="required" type="xs:string"/>
                  <xs:attribute name="unit" use="required" type="xs:string"/>
                </xs:extension>
              </xs:simpleContent>
            </xs:complexType>
          </xs:element>

        </xs:choice>
      </xs:sequence>
    </xs:complexType>
  </xs:element>

</xs:schema>
    \end{verbatim}

    Datagrama JSON (ejemplo):
    \begin{verbatim}
    {
        "server": "<id de servidor>",
        "measurements": [
            {
                "magnitude": "<magnitud>",
                "unit": "<unidad>",
                "value": "<valor>"
            }
        ]
    }
    \end{verbatim}

    Comandos:
    \begin{verbatim}
<?xml version="1.0"?>
<xs:schema xmlns:xs="http://www.w3.org/2001/XMLSchema">

  <xs:element name="control">
    <xs:complexType>
      <xs:choice minOccurs="0" maxOccurs="unbounded">
        <xs:element name="format">
            <xs:simpleType>
                <xs:restriction base="xs:string">
                    <xs:enumeration value="json"/>
                    <xs:enumeration value="xml"/>
                </xs:restriction>
            </xs:simpleType>
        </xs:element>
        <xs:element name="frequency" type="xs:integer"/>
        <xs:element name="unit">
          <xs:complexType>
            <xs:simpleContent>
              <xs:extension base="xs:string">
                <xs:attribute name="magnitude" use="required" type="xs:string"/>
              </xs:extension>
            </xs:simpleContent>
          </xs:complexType>
        </xs:element>
        <xs:element name="enable" type="xs:boolean"/>
      </xs:choice>
    </xs:complexType>
  </xs:element>

</xs:schema>
    \end{verbatim}

    \subsubsection{Servidor}

    Aplicación en Java con dos threads, uno envía periodicamente datagramas,
    y el otro está a la espera de comandos. Se compone de una sola clase que implementa
    Runnable para lanzar el thread que escucha comandos, mientras que el main thread
    es el que se dedica a enviar datagramas.

    \subsubsection{Cliente}

    Aplicación en Java con un solo thread. Según se ejecute, se queda escuchando datagramas
    en una interfaz, o envía un comando y termina. También se compone de una sola clase.

    \section{Ejecución}

    \subsection{Compilación}
    
    \begin{verbatim}
    make
    \end{verbatim}
    
    \subsection{Lanzamiento}

    \subsubsection{Servidor}

    \begin{verbatim}
    java -jar build/server.java [interfaz] [id de servidor] [weather | quality]
    \end{verbatim}

    \begin{itemize}
    \item interfaz: interfaz donde enviar tráfico multicast
    \item id de servidor: campo que identifica al servidor que manda los datagramas
    \item weather \textbar quality: dos tipos de sensores en el servidor
    \end{itemize}

    Necesario especificar interfaz

    \subsubsection{Cliente}

    \begin{verbatim}
    java -jar build/client.java [-i <interfaz>] <listen | cmd <servidor>
     <comando> <argumento> [magnitud]]
    \end{verbatim}

    \begin{itemize}
    \item interfaz: interfaz donde recibir tráfico multicast
    \item listen \textbar cmd: escuchar datos de sensores, o enviar un comando
    \item servidor: servidor al que mandarle comando de control
    \item comando: comando que enviar
    \begin{itemize}
        \item format \textless json \textbar xml \textgreater: formato de datagramas de datos
        \item frequency \textless f\textgreater: frecuencia en Hz de muestreo
        \item unit \textless unidad\textgreater \textless magnitud\textgreater: unidad para la magnitud
        \item enable \textless true \textbar false \textgreater: encender o apagar envio de datos
    \end{itemize}
    \item argumento: argumento del comando enviado
    \item magnitud: el comando
    \end{itemize}

    Necesario especificar interfaz (\texttt{lo} por defecto)

    El cliente puede recibir datagramas en cualquiera de los dos formatos, XML o JSON,
    de manera transparente.

    \section{Autoría}

    Declaro que he confeccionado la totalidad de esta práctica.
    No se ha usado IA tales como LLMs de ninguna manera para realizar esta práctica.

\end{document}

